%-------------------------
% Resume in Latex
% Author : Jake Gutierrez
% Based off of: https://github.com/sb2nov/resume
% License : MIT
%------------------------

\documentclass[letterpaper,11pt]{article}

\usepackage{latexsym}
\usepackage[empty]{fullpage}
\usepackage{titlesec}
\usepackage{marvosym}
\usepackage[usenames,dvipsnames]{color}
\usepackage{verbatim}
\usepackage{enumitem}
\usepackage[hidelinks]{hyperref}
\usepackage{fancyhdr}
\usepackage[english]{babel}
\usepackage{tabularx}
\usepackage{multicol}
\setlength{\multicolsep}{-3.0pt}
\setlength{\columnsep}{-1pt}
\input{glyphtounicode}
\usepackage{graphicx}
\usepackage[absolute,overlay]{textpos}


%----------FONT OPTIONS----------
% sans-serif
% \usepackage[sfdefault]{FiraSans}
% \usepackage[sfdefault]{roboto}
% \usepackage[sfdefault]{noto-sans}
% \usepackage[default]{sourcesanspro}

% serif
% \usepackage{CormorantGaramond}
% \usepackage{charter}


\pagestyle{fancy}
\fancyhf{} % clear all header and footer fields
\fancyfoot{}
\renewcommand{\headrulewidth}{0pt}
\renewcommand{\footrulewidth}{0pt}

% Adjust margins
\addtolength{\oddsidemargin}{-0.5in}
\addtolength{\evensidemargin}{-0.5in}
\addtolength{\textwidth}{1in}
\addtolength{\topmargin}{-.5in}
\addtolength{\textheight}{1.0in}

\urlstyle{same}

\raggedbottom
\raggedright
\setlength{\tabcolsep}{0in}

% Sections formatting
\titleformat{\section}{
  \vspace{-5pt}\scshape\raggedright\large
}{}{0em}{}[\color{black}\titlerule \vspace{-5pt}]

% Ensure that generate pdf is machine readable/ATS parsable
\pdfgentounicode=1

%-------------------------
% Custom commands
\newcommand{\resumeItem}[1]{
  \item\small{
    {#1 \vspace{-2pt}}
  }
}

\newcommand{\resumeSubheading}[4]{
  \vspace{-2pt}\item
    \begin{tabular*}{0.97\textwidth}[t]{l@{\extracolsep{\fill}}r}
      \textbf{#1} & #2 \\
      \textit{\small#3} & \textit{\small #4} \\
    \end{tabular*}\vspace{-7pt}
}

\newcommand{\resumeSubSubheading}[2]{
    \item
    \begin{tabular*}{0.97\textwidth}{l@{\extracolsep{\fill}}r}
      \textit{\small#1} & \textit{\small #2} \\
    \end{tabular*}\vspace{-7pt}
}

\newcommand{\resumeProjectHeading}[2]{
    \vspace{-2pt}\item
    \begin{tabular*}{0.97\textwidth}{l@{\extracolsep{\fill}}r}
      \small#1 & #2 \\
    \end{tabular*}\vspace{-7pt}
}

\newcommand{\resumeSubItem}[1]{\resumeItem{#1}\vspace{-4pt}}

\renewcommand\labelitemii{$\vcenter{\hbox{\tiny$\bullet$}}$}

\newcommand{\resumeSubHeadingListStart}{\begin{itemize}[leftmargin=0.15in, label={}]}
\newcommand{\resumeSubHeadingListEnd}{\end{itemize}}
\newcommand{\resumeItemListStart}{\begin{itemize}}
\newcommand{\resumeItemListEnd}{\end{itemize}\vspace{-5pt}}

%-------------------------------------------
%%%%%%  RESUME STARTS HERE  %%%%%%%%%%%%%%%%%%%%%%%%%%%%


\begin{document}

%----------HEADING----------
\begin{tabular*}{\textwidth}{l@{\extracolsep{\fill}}l r}
\begin{tabular}[t]{l}
\textbf{\Huge \scshape Genova Mongalo}\\
\small 913-260-8912 $|$ 
\href{mailto:Genova.Mongalo@gmail.com}{\underline{Genova.Mongalo@gmail.com}} $|$ 
\href{https://linkedin.com/in/GMongalo}{\underline{linkedin.com/in/GMongalo}} $|$
\href{https://gmango.dev}{\underline{gmango.dev}}
\end{tabular}
\end{tabular*}
% The QR code goes to GitHub Pages, the fallback that outlives a lapsed domain.
% Same site as gmango.dev, so an ATS reading the text still gets a working link.
\begin{textblock*}{30mm}(178.2mm,8mm)
  \href{https://ShankyShako.GitHub.io}{\includegraphics[width=25mm, trim={4cm 4cm 4cm 4cm},clip]{qr-code.png}}
\end{textblock*}

%-----------EDUCATION-----------
\section{Education}
  \resumeSubHeadingListStart
    \resumeSubheading
      {Georgia Institute of Technology}{Atlanta, GA}
      {Master of Science in Computer Science, emphasis in Artificial Intelligence}{Expected 2027}
    \resumeSubheading
      {University of Missouri-Kansas City (UMKC)}{Kansas City, MO}
      {Bachelor of Science in Computer Science, summa cum laude}{Aug. 2021 -- Dec. 2024}
  \resumeSubHeadingListEnd


%-----------EXPERIENCE-----------
\section{Experience}
  \resumeSubHeadingListStart
    \resumeSubheading
      {Founding AI/ML Product Engineer}{September 2025 -- Present}
      {Geometry Health and Wellness}{Kansas City, MO}
      \resumeItemListStart
        \resumeItem{Ported clinically validated mobility assessments from MATLAB prototypes to a production iOS app, within 0.56\% of the clinical reference on a 15-biomarker postural-sway test and exact on 10-meter-walk gait speed.}
        \resumeItem{Built the on-device biomarker pipeline in Swift: 100 Hz CoreMotion capture, gravity removal, zero-phase Butterworth band-pass filtering, double integration, and PCA into 15 biomarkers with quality flags.}
        \resumeItem{Built an end-to-end clinical LLM assistant with patient-context retrieval and PubMed evidence via Entrez, self-hosted on Ollama (DeepSeek-V2 16B) and streaming over Flask REST + WebSocket in under 1 s.}
        % \resumeItem{Built a speech assessment pipeline pairing on-device Apple Speech for live partials with a Whisper backend and LLM transcript normalization, with a patient-review panel, to catch the communication deficits that affect one in three stroke survivors.}
        % \resumeItem{Built the in-house iOS data-collection app for motion, survey, and fall/emergency data; each trial exports as timestamped JSON into the labeled datasets the models are validated against.}
        \resumeItem{Led the product from MVP to production as founding engineer, owning the architecture and sprint cycles and keeping the clinical, engineering, and product sides in step.}
      \resumeItemListEnd
    \resumeSubheading
      {Federal AI/ML Engineer Contractor}{May 2025 -- August 2025}
      {AFRL Sensors Directorate Contracting Program, University of Dayton}{Dayton, OH}
      \resumeItemListStart
        \resumeItem{Built generative AI models that let a CNN classifier recognize real objects after training only on generated data, carried over by transfer learning, for AFRL/DoD.}
        \resumeItem{Devised a game-theoretic GAN training scheme, run on Slurm-parallelized compute; training on its output raised downstream classifier accuracy 50 percentage points over the baseline.}
        % \resumeItem{Adapted state-of-the-art PyTorch classifiers to imaging modalities they were not designed for.}
      \resumeItemListEnd
    \resumeSubheading
      {AI/ML Engineer Research Intern}{June 2024 -- Dec. 2024}
      {NSF REU AI-Empowered Cybersecurity - University of Missouri, Kansas City}{Kansas City, MO}
      \resumeItemListStart
        \resumeItem{Built a transformer-based model that detects ransomware in Industrial Control Systems at 99\% binary and 91\% family classification accuracy; presented it at IEEE Big Data 2024 as solo author, now accepted for publication.}
      \resumeItemListEnd

    % \resumeSubheading
    %   {Research Intern}{May 2023 -- July 2023}
    %   {NSF REU in Consumer Networking - University of Missouri, Columbia}{Columbia, MO}
    %   \resumeItemListStart
    %     \resumeItem{Designed a transformer model that predicts cybersickness severity on a 1-10 scale with 85\% accuracy.}
    %   \resumeItemListEnd

    % \resumeSubheading
    %   {Research Intern}{Dec. 2021 -- April 2022}
    %   {NASA Missouri Space Grant Consortium}{}
    %   \resumeItemListStart
    %     \resumeItem{Engineered an AI-driven simulation of space-bound transmissions for a NASA-funded study of signal propagation, finding which frequency bands hold up under which conditions.}
    %     \resumeItem{Developed an interactive interface that shows transmission strength and recommends a frequency for the environment you give it.}
    %   \resumeItemListEnd
  \resumeSubHeadingListEnd


%-----------PROJECTS-----------
\section{Projects}
    \resumeSubHeadingListStart
      \resumeProjectHeading
          {\textbf{LegacyScribe} $|$ \emph{Swift, SwiftUI, MLX, Speech Recognition, LLMs}}{}
          \resumeItemListStart
            \resumeItem{Built an on-device macOS dictation app I use daily (a SwiftUI fork of Hex) with swappable speech engines: 5.3\% word error rate on Cohere Transcribe/MLX vs 8.8\% on Parakeet, 10 s of audio in about 0.33 s.}
            % \resumeItem{Added an LLM pass that cleans up a transcript or reshapes a spoken ramble into a structured prompt for a coding agent, and pastes the raw transcript if every model endpoint fails.}
            % Optional (resume.ts: scribe-ondevice, scribe-private, scribe-practice):
            % \resumeItem{Runs transcript rewriting on-device with Qwen3 4B (4-bit) in 0.4--1.5 s for dictations of up to 55 words; the live caption shows its first word a median 0.74 s after speech starts.}
            % \resumeItem{Built a private mode that opens zero network connections, verified by a test that samples the app's sockets every 10 ms during a real dictation.}
            % \resumeItem{Added a read-aloud practice mode built from 280 sentences covering all 39 English sounds and 881 sound pairs, plus vocabulary suggestions drawn from what the user actually sent after editing.}
          \resumeItemListEnd
      \resumeProjectHeading
          {\textbf{Care Beyond (checkpointkc.com)} $|$ \emph{React, Leaflet, Node.js, Express}}{}
          \resumeItemListStart
            \resumeItem{Built and deployed a live map of food, shelter, clothing, and medical resources from local organizations, searchable by address or ZIP within a radius; only verified organizations post, signing in with a one-time email code.}
          \resumeItemListEnd
      \resumeProjectHeading
          {\textbf{Dysarthria Speech Understanding} $|$ \emph{PyTorch, Whisper, wav2vec2}}{}
          \resumeItemListStart
            \resumeItem{Built the ASR evaluation harness and Whisper multi-task grid for dysarthric speech, run leave-one-speaker-out across 43 speakers: 10.44 MAE on intelligibility and 0.809 quadratic weighted kappa on severity.}
            % Optional (resume.ts: dys-findings):
            % \resumeItem{Measured transcription degrading roughly thirteen-fold from control to impaired speakers on identical prompts, with Whisper and wav2vec2 failing in opposite directions: fluent but wrong versus visibly broken.}
          \resumeItemListEnd
      % \resumeProjectHeading
      %     {\textbf{AWS DeepRacer} $|$ \emph{Reinforcement Learning, TD3, Docker}}{}
      %     \resumeItemListStart
      %       \resumeItem{Hand-shaped the reward for a TD3 policy driving a simulated 1/18-scale car: center-line distance bands, heading toward upcoming waypoints, a speed-versus-steering trade-off, and penalties for leaving the track.}
      %     \resumeItemListEnd
      % Optional projects (resume.ts: prj-site, prj-cogload):
      % \resumeProjectHeading
      %     {\textbf{gmango.dev} $|$ \emph{React, TypeScript, LLMs}}{}
      %     \resumeItemListStart
      %       \resumeItem{Built gmango.dev's r\'esum\'e tailoring: in-browser matching re-ranks a bullet library on every keystroke, a typesetter keeps the page to one sheet, a language model rewords, and a rule-based checker rejects any rewrite that drops a number, adds a tool, or claims more.}
      %       \resumeItem{Built the site's chat assistant, grounded in a knowledge file generated from the site's own data at build time; it can only attach links from a generated menu, so it never sends a visitor to a page that does not exist.}
      %     \resumeItemListEnd
      % \resumeProjectHeading
      %     {\textbf{Cognitive-Load-Aware Conversational Design} $|$ \emph{HCI, Cognitive Science}}{}
      %     \resumeItemListStart
      %       \resumeItem{Analyzed a deployed clinical conversational agent's interface and dialogue against cognitive load theory, Baddeley's working-memory model, dual-process theory, and Levelt's speech-production model (no human subjects).}
      %     \resumeItemListEnd
    \resumeSubHeadingListEnd


%-----------PUBLICATIONS-----------
\section{Publications}
    \begin{itemize}[leftmargin=0.15in, label={}, itemsep=0pt]
      \item\small{Mongalo, G. (2026). An Explainable Machine Learning Approach and Provenance Registry for Real-Time Ransomware Detection in ICS. IEEE Big Data 2026, accepted, to appear Dec. 2026 (sole author).}
      \item\small{Kundu, R. K., Denton, M., Mongalo, G., et al. (2025). Securing Virtual Reality Experiences: Unveiling and Tackling Cybersickness Attacks With Explainable AI. IEEE Trans. Dependable Secure Comput., 22(6).}
    \end{itemize}


%-----------PROGRAMMING SKILLS-----------
\section{Technical Skills}
 \begin{itemize}[leftmargin=0.15in, label={}]
    \item \small
     \textbf{Languages}{: Python, Swift, C / C++, SQL (Postgres), JavaScript, Java / C\#}\\
     \textbf{Machine Learning}{: PyTorch, scikit-learn, Hugging Face Transformers, Whisper, pandas, NumPy}\\
     \textbf{Frameworks}{: SwiftUI, Flask, FastAPI, Socket.IO, React, Node.js}\\
     \textbf{Developer Tools}{: Git, Docker, Slurm, AWS, Google Cloud Platform, Xcode}
 \end{itemize}


%-------------------------------------------
\end{document}
